\chapter{Connections}
\label{ch:connections}

This chapter lists what every signal of every chip connects to. A net is one wire that joins the pins named in the row.
A pin that is not in any table is a power pin, or has no connection that this manual can state.

\section{The Shared Bus}
\label{sec:shared-bus}

The three chips share the signals of Table~\ref{tab:shared-nets}. Pin numbers are those of the packages: the CPU is a
64-pin DIP and SIRCADA a 52-pin PLCC (\hwmanual{}, \videomanual{}). The package of SIRCAS is not yet specified, so it
has no pin numbers.

\begin{table}[H]
	\centering
	\small
	\begin{tabular}{lllll}
		\toprule
		\textbf{Net}           & \textbf{CPU} & \textbf{SIRCADA} & \textbf{SIRCAS}        & \textbf{Other}                    \\
		\midrule
		A0--A7                 & 1--8         & 8--15            & A0--A7                 & Memories, cartridge               \\
		A8--A15                & 10--17       & 16--23           & A8--A15                & Memories, cartridge               \\
		A16--A23               & 18--25       & ---              & ---                    & Select logic, memories, cartridge \\
		D0--D15                & 27--42       & 26--41           & D0--D15                & Memories, cartridge               \\
		\activelow{BAS}        & 47 (out)     & 5 (in)           & \activelow{BAS} (in)   & Select logic                      \\
		BRW                    & 48 (out)     & 6 (in)           & BRW (in)               & Memories, cartridge               \\
		\activelow{BACK}       & 49 (in)      & 7 (out)          & \activelow{BACK} (out) & Memories, cartridge; pull-up      \\
		\activelow{CS} (video) & ---          & 4                & ---                    & Select logic, Y3                  \\
		\activelow{CS} (audio) & ---          & ---              & \activelow{CS}         & Select logic, Y4                  \\
		CLKI                   & 45           & 2                & CLKI                   & Clock circuit                     \\
		\bottomrule
	\end{tabular}
	\caption{Nets Shared by the Chips}
	\label{tab:shared-nets}
\end{table}

\noindent The data and address bits connect bit for bit: D0 of the CPU to D0 of each device, and so on.
\activelow{BACK} is open drain at the devices, so it is a wired net, pulled up by a resistor whose value is not yet
specified. The memories and the cartridge may take part in it too. The \activelow{CS} of each device is a separate net from the
select logic.

\section{Interrupts and Reset}

\begin{table}[H]
	\centering
	\small
	\begin{tabular}{lllll}
		\toprule
		\textbf{Net}     & \textbf{CPU} & \textbf{SIRCADA}       & \textbf{SIRCAS}        & \textbf{Notes}             \\
		\midrule
		\activelow{NMI}  & 62 (in)      & 43 (out)               & ---                    & Pull-up; vertical sync     \\
		\activelow{IRQ4} & 61 (in)      & 44 (out)               & ---                    & Pull-up; line compare      \\
		\activelow{IRQ3} & 60 (in)      & ---                    & \activelow{IRQ3} (out) & Pull-up; sequence complete \\
		\activelow{IRQ2} & 59 (in)      & ---                    & ---                    & Not yet specified          \\
		\activelow{IRQ1} & 58 (in)      & ---                    & ---                    & Not yet specified          \\
		\activelow{RSTO} & 64 (out)     & 3 (\activelow{RST} in) & \activelow{RST} (in)   & One net to both devices    \\
		\activelow{RSTI} & 63 (in)      & ---                    & ---                    & Reset circuit              \\
		\bottomrule
	\end{tabular}
	\caption{Interrupt and Reset Nets}
	\label{tab:irq-reset-nets}
\end{table}

\section{Signals the Devices Do Not Use}

The CPU drives and receives more signals than the other two chips use. What SIRCUS does with them is as follows.

\begin{table}[H]
	\centering
	\small
	\begin{tabular}{lll}
		\toprule
		\textbf{CPU signal} & \textbf{Pin} & \textbf{Connection}                                             \\
		\midrule
		\activelow{BERR}    & 50 (in)      & Driven by the system for an unmapped segment; not yet specified \\
		\activelow{BPER}    & 51 (in)      & Driven by a memory controller, if any; not yet specified        \\
		BAT0--BAT2          & 52--54 (out) & Memory controller, if any; not yet specified                    \\
		PROT                & 55 (out)     & Memory controller, if any; not yet specified                    \\
		SYNC                & 46 (out)     & Not yet specified                                               \\
		\activelow{HALT}    & 56 (in)      & Not yet specified                                               \\
		\activelow{TRCE}    & 57 (in)      & Not yet specified                                               \\
		\bottomrule
	\end{tabular}
	\caption{CPU Signals Not Shared with the Other Chips}
	\label{tab:unused-cpu}
\end{table}

\section{The Select Logic}
\label{sec:select-logic}

One 74LS138 decoder makes every chip select (Section~\ref{sec:selecting}).

\begin{table}[H]
	\centering
	\small
	\begin{tabular}{lllp{5.6cm}}
		\toprule
		\textbf{74LS138 pin} & \textbf{Name} & \textbf{Connection}            & \textbf{Notes}                                                \\
		\midrule
		1                    & A             & CPU A21 (pin 23)               &                                                               \\
		2                    & B             & CPU A22 (pin 24)               &                                                               \\
		3                    & C             & CPU A23 (pin 25)               &                                                               \\
		4                    & G2A           & \activelow{BAS}                & Enable, active low                                            \\
		5                    & G2B           & Ground                         & Enable, active low                                            \\
		6                    & G1            & $+5$ V                         & Enable, active high                                           \\
		15                   & Y0            & Working memory                 & Gated by the boot overlay; see Section~\ref{sec:boot-overlay} \\
		14                   & Y1            & System ROM                     &                                                               \\
		13                   & Y2            & Cartridge                      &                                                               \\
		12                   & Y3            & SIRCADA \activelow{CS} (pin 4) &                                                               \\
		11                   & Y4            & SIRCAS \activelow{CS}          &                                                               \\
		10                   & Y5            & Serial I/O                     &                                                               \\
		9, 7                 & Y6, Y7        & Not yet specified              & Could drive \activelow{BERR}                                  \\
		8, 16                & GND, VCC      & Ground, $+5$ V                 &                                                               \\
		\bottomrule
	\end{tabular}
	\caption{The Select Decoder}
	\label{tab:select-decoder}
\end{table}

\noindent The gating of the working memory's select and of the cartridge's by the boot overlay, and the way each memory
produces \activelow{BACK}, are not yet specified.

\section{The Serial Device}

The serial device is a basic serial port. Its purpose is to link a console to another console for two-player games, to
a modem, or to a debugger. The part, its signals, its speed and its connector are not yet specified. It takes its
\activelow{CS} from output Y5 of the decoder and the signals of the shared bus (Table~\ref{tab:shared-nets}), and it
needs an interrupt (Chapter~\ref{ch:interrupts}).

\section{SIRCAS and the Audio Memory}
\label{sec:audio-memory-wiring}

The audio memory is a static memory of 32,768 words of 16 bits, which may be two memories of 32,768 by 8 bits side by
side. It connects to SIRCAS and to nothing else.

\begin{table}[H]
	\centering
	\begin{tabular}{lll}
		\toprule
		\textbf{SIRCAS} & \textbf{Audio memory}                  & \textbf{Notes} \\
		\midrule
		MA0--MA14       & A0--A14 of both memories               & Word address   \\
		MD0--MD7        & D0--D7 of the memory for the low byte  & Three-state    \\
		MD8--MD15       & D0--D7 of the memory for the high byte & Three-state    \\
		\activelow{MCE} & Chip enable of both                    &                \\
		\activelow{MOE} & Output enable of both                  &                \\
		\activelow{MWE} & Write enable of both                   &                \\
		\bottomrule
	\end{tabular}
	\caption{Audio Memory Connections}
	\label{tab:audio-memory-nets}
\end{table}

A transaction occupies two \texttt{CLKI} cycles, 186 ns in an NTSC system and 188 ns in a PAL system, so the memory
must complete an access within that time less the set-up and hold times of the interface. Those times, and the part
chosen for the memory, are not yet specified (\audiomanual{}, package chapter). The mapping of data bits to the
memories follows the byte order of the audio memory (\audiomanual{}, memory map chapter).

\section{SIRCAS and the Converter}
\label{sec:converter}

SIRCAS sends its output to the converter as a serial stream. The converter is the Philips TDA1541A, a stereo 16-bit
converter with current outputs, in a 28-pin DIP. It is used in its time-multiplexed two's-complement mode, which takes
the word select, the bit clock and one data line.

\begin{table}[H]
	\centering
	\small
	\begin{tabular}{llp{6.4cm}}
		\toprule
		\textbf{SIRCAS} & \textbf{TDA1541A}             & \textbf{Notes}                                                                 \\
		\midrule
		BCLK            & BCK, pin 2                    & 10 and 11 \texttt{CLKI} cycles alternately: 1.023 MHz mean NTSC, 1.013 MHz PAL \\
		WS              & LE/WS, pin 1                  & Low for left, high for right; the sample rate                                  \\
		SD              & DATA L/DATA, pin 3            & Signed 16-bit samples, most significant bit first                              \\
		SCK             & ---                           & Not connected: the converter has no system clock input                         \\
		---             & DATA R, pin 4                 & Not used in this mode                                                          \\
		---             & \activelow{OB}/TWC, pin 27    & Tied to +5 V: time-multiplexed, two's complement                               \\
		---             & COSC, pins 16 and 17          & A 470 pF capacitor between them: the converter's own oscillator                \\
		---             & DECOU, pins 7 to 13, 18 to 24 & A 100 nF capacitor from each to ground                                         \\
		---             & AOL, pin 25; AOR, pin 6       & The current outputs, one to each output amplifier                              \\
		\bottomrule
	\end{tabular}
	\caption{Converter Connections}
	\label{tab:converter-nets}
\end{table}

\subsection{Converter Timing}

Table~\ref{tab:converter-timing} sets what the converter requires of its inputs against what SIRCAS gives.

\begin{table}[H]
	\centering
	\small
	\begin{tabular}{p{4.6cm}p{2.9cm}p{4.6cm}}
		\toprule
		\textbf{Parameter}                                    & \textbf{Converter needs} & \textbf{SIRCAS gives}                                                     \\
		\midrule
		Bit clock cycle                                       & 156 ns or more           & 931 or 1,024 ns NTSC, 940 or 1,034 ns PAL (10 or 11 \texttt{CLKI} cycles) \\
		Bit clock high time, low time                         & 46 ns or more each       & At least 93 ns each, whatever the split of the 10 or 11 cycles            \\
		Word select rate                                      & 200 kHz or less          & 31.96 kHz NTSC, 31.67 kHz PAL                                             \\
		Rise and fall times                                   & 32 ns or less            & Not yet specified                                                         \\
		Data and word select set-up before \texttt{BCK} rises & 32 ns or more            & Not yet specified                                                         \\
		Data and word select hold after \texttt{BCK} rises    & 0 ns or more             & Not yet specified                                                         \\
		\bottomrule
	\end{tabular}
	\caption{Converter Input Timing}
	\label{tab:converter-timing}
\end{table}

\noindent The bit clock is far inside the converter's limits. Whether the data edges meet the set-up time depends on which
edge of \texttt{BCLK} SIRCAS changes them on, which its manual leaves open; a change on the falling edge leaves the low time of
the bit clock, at least one \texttt{CLKI} cycle, 93 ns, whatever the duty cycle. The converter accepts words of up to 16
bits, most significant bit first, and SIRCAS sends exactly 16 in a slot of 16 bit clocks, so nothing follows the sixteenth.
As in the standard interface, the last bit of each word is sent in the first bit clock of the next slot, after the word select
has changed. The bit clock periods alternate between 10 and 11 cycles in a pattern that repeats every sample, so the edges
fall at the same cycle of the sample each time.

\subsection{The Output Amplifiers}

Each output of the converter is a current that sinks up to 4.0 mA at full scale (3.4 to 4.6 mA), and must be held at
virtual earth within $\pm 25$~mV. A reference circuit for the part converts it to a voltage with an operational
amplifier for each channel, the inverting input at the converter output, the other input at analogue ground, and a
1.8~k$\Omega$ resistor with a 2.2~nF capacitor in parallel in the feedback. The full-scale output of that circuit is
4.0~mA times 1.8~k$\Omega$, 7.2~V. The capacitor and resistor give a corner at 40~kHz, which is above half the sample
rate here, so the circuit does not by itself remove the images of the signal around 32~kHz; a steeper filter follows.
The amplifier, its supplies, the filter and the output connector are not yet specified.

\section{SIRCADA and the Display}

SIRCADA's video output leaves the system as an analogue colour signal and a composite sync. The video memory is inside
the device, so it has no memory interface.

\begin{table}[H]
	\centering
	\begin{tabular}{lll}
		\toprule
		\textbf{SIRCADA}  & \textbf{Pin} & \textbf{Connection}                               \\
		\midrule
		R, G, B           & 48--50       & Display connector, through the output circuit     \\
		\activelow{CSYNC} & 45           & Display connector                                 \\
		IREF              & 47           & A reference that sets the full-scale colour level \\
		AVCC, AGND        & 46, 51       & The analogue supply and ground                    \\
		\bottomrule
	\end{tabular}
	\caption{Display Connections}
	\label{tab:display-nets}
\end{table}

\noindent The output circuit, the reference, the connector and the levels at it are not yet specified.
